\chapter{Итоговая заморозка родительской архитектуры}

Пятый том был открыт не для получения ещё одного частного потенциала, а для
проверки более сильного утверждения: существует ли в пределах принятой
программы один родительский объект, который до феноменологии фиксирует
носитель, состояние, меру, ориентацию, физический фактор и относительные
нормировки секторов. Настоящая глава сводит все машинные результаты Тома V
и применяет условие ранней остановки из первой главы.

\section{Область утверждения}

Итог относится только к трём заранее объявленным классам:
\begin{enumerate}
 \item совместной вариации состояния и носителя;
 \item единому граничному источнику;
 \item конечной, скрученной или минимально производной геометрии в
 замороженном бюджете сложности.
\end{enumerate}
Он не является запретом всех возможных квантово-гравитационных,
модулярных или симплектических теорий. Он устанавливает исчерпание
конкретного меню Тома V при требованиях положительной меры, вещественной
структуры KO6, сохранения физической группы и отсутствия внешних
межсекторных весов.

\section{Совместная вариация состояния и носителя}

Для фиксированного носителя вариация Гиббса правильно нормирует состояние.
Но родительская мера на метриках и топологиях не получена; две допустимые
семантики спектрального функционала меняют порядок предпочтительных
носителей; конечные коэффициенты кривизны также способны обращать этот
порядок. Полная мера полей, духов и массивных векторов не задана, поэтому
совместный гессиан не определён.

Попытка восстановить геометрию из корреляционного объекта через треугольник
редукций также не замкнулась. Слабые корреляционные данные недоопределяют
геометрию, а включение полной спектральной тройки в исходный объект делает
обратное отображение тавтологическим.

\begin{equation}
 \boxed{\text{нормировка состояния --- да, родительская мера --- нет}.}
\end{equation}

Следовательно, текущая реализация первого класса закрыта до вычисления
спектральных сумм и наблюдаемых.

\section{Архитектура граничного источника}

Граничная ветвь сохраняет несколько точных локальных модулей:
вильсоновскую пару коэффициентов, поколенческую ось, условный конденсат и
одну моду Майораны. Однако они действуют в нескольких неэквивалентных
блоках. Симметрийный след на их прямой сумме оставляет свободные центральные
веса, а простое условие Гаусса не выбирает нужный многокомпонентный сектор
заряда.

Операторнозначный проектор или мнимый когерентный источник могли бы выбрать
требуемую ветвь, но являлись бы дополнительными данными. Поэтому один
граничный след не выводит одновременно заряд, конденсат, тетраэдрическую
ось, ядро Майораны и вторую независимую нормировку.

\begin{equation}
 \boxed{\text{локальные модули --- частично да, единый источник --- нет}.}
\end{equation}

\section{Конечная, скрученная и производная геометрия}

Эта ветвь была исследована наиболее подробно.
\begin{enumerate}
 \item Ограниченный перебор конечных графов выполнен, но висячий лист не
 выбирает поколенческую матрицу.
 \item Активный прямоугольник алгебры поколений проходит условие первого
 порядка и содержит чувствительный к циклу квартичный член, но его глобальный
 минимум обнуляет триплетные связующие поля.
 \item Обычный спектральный функционал видит \(A+B\), но не восстанавливает
 ориентированную разность \(A-B\).
 \item Действие высоты--Ходжа точно воспроизводит нужный квадрат и правильную
 нормировку, но данные KO6 допускают как согласованную цепь, так и сток.
 \item Селективное скручивание даёт допустимое вещественное представление и
 скрученное условие первого порядка, но не создаёт нового блока кривизны и
 не исправляет знак обычного следа.
 \item Точный скрученный след аннулирует переставляемые центральные
 идемпотенты, а реализующий обмен вес неизбежно вырожден.
 \item Предпроективное соотношение точно даёт
 \(XX^\dagger-Y^\dagger Y\), но его формальная независимость от ориентации
 несовместима с положительной инволюцией. Для физической
 \(\mathfrak{so}(3)\) симметрическая матрица имеет нулевое следовое
 спаривание.
\end{enumerate}

Производный комплекс после задания соотношения нильпотентен, но это
разрешение уже выбранного отображения момента, а не его происхождение из
текущей спектральной геометрии.

\begin{equation}
 \boxed{\text{ориентированное тождество --- да, его родительский выбор --- нет}.}
\end{equation}

\section{Положительный математический остаток}

Отрицательный итог архитектуры не обнуляет результаты Тома V. Сохраняются:
\begin{enumerate}
 \item точная нормировка состояния при фиксированном носителе;
 \item конечная классификация малых графов внутри объявленного бюджета;
 \item явный прямоугольник первого порядка и скалярный коммутант семейной
 алгебры;
 \item теорема спектральной слепоты к ориентированной разности матриц Грама;
 \item точное тождество высоты--Ходжа и нормировка полного следа;
 \item явное вещественное скрученное представление KO6;
 \item запрет точного положительного следа для центрального обмена;
 \item точное предпроективное средневершинное соотношение;
 \item запрет совместить его независимость от ориентации с положительной
 инволюцией;
 \item препятствие строгого отображения момента группы \(SO(3)\).
\end{enumerate}

Это самостоятельный математический результат уровня классификации и
запретительных теорем. Он не является физическим замыканием.

\section{Проверка критерия завершённости}

\begin{center}\small
\begin{tabular}{>{\raggedright\arraybackslash}p{7.2cm}
                  >{\raggedright\arraybackslash}p{2.2cm}
                  >{\raggedright\arraybackslash}p{3.4cm}}
\toprule
Требование & Статус & Основание\\
\midrule
Выбор направления до феноменологии & выполнено & наблюдаемые входы не использованы\\
Один заранее фиксированный функционал & не выполнено & мера и ориентация остаются внешними\\
Полная мера вспомогательных секторов & не выполнено & поля, духи и веса не объединены\\
Физический фактор БВ/БРСТ и полный гессиан & не выполнено & положительного кандидата нет\\
Один след для двух независимых секторов & не выполнено & относительные веса свободны\\
Пороговый ход без целевых масс & не начинался & архитектура не допущена\\
Две независимые слепые проверки & не начинались & феноменология запрещена\\
\bottomrule
\end{tabular}
\end{center}

\begin{theorem}[Ранняя остановка Тома V]
В пределах трёх объявленных классов, замороженного конечного бюджета,
положительной гильбертовой инволюции, текущей вещественной группы и запрета
на внешние межсекторные веса не существует прошедшей проверку V.0
родительской архитектуры. Поэтому переход к подгонке масс, смешиваний,
связей, порогов или абсолютного масштаба не разрешён.
\end{theorem}

Это утверждение является исчерпанием проверенного меню, а не универсальным
запретом будущей теории.

\section{Условия повторного открытия}

Продолжение возможно только после явного объявления нового родительского
объекта. Он обязан:
\begin{enumerate}
 \item вывести ориентацию или поляризацию до записи квадрата отображения
 момента;
 \item вывести физическую вещественную форму и любое расширение
 калибровочной группы;
 \item дать положительную меру и один след как минимум для двух независимых
 секторов;
 \item определить полный фактор БВ/БРСТ и физический гессиан до
 феноменологии;
 \item отличаться от уже закрытых реализаций новым проверяемым
 математическим содержанием.
\end{enumerate}

К возможным, но не открытым здесь классам относятся комплексный
симплектический родитель с выведенным положительным физическим сектором,
модулярная алгебра типа III с конкретной конечной редукцией и неприводимое
граничное гильбертово пространство с выведенными межблочными связями.
Каждый из них является новой аксиоматикой, а не поправкой к текущему тому.

\section{Итоговый статус}

Оценка \(R_{\rm sci}=4/10\) наследуется из Четвёртого тома и не
пересчитывается без отдельной шкальной проверки. Число математически
замкнутых родительских архитектур, межсекторных замыканий и физических
замыканий равно нулю:
\begin{equation}
 N_{\rm parent}=N_{\rm intersector}=N_{\rm physical}=0.
\end{equation}

\begin{protocol}
\begin{enumerate}
 \item процедурная проверка до феноменологии: \textbf{выполнена};
 \item положительные математические тождества и запреты:
 \textbf{сохранены};
 \item критерий завершённости Тома V: \textbf{не выполнен};
 \item текущее меню родительских архитектур: \textbf{исчерпано};
 \item условие ранней остановки: \textbf{сработало};
 \item феноменологическая подгонка: \textbf{запрещена};
 \item физическое замыкание: \textbf{не достигнуто};
 \item Том V: \textbf{заморожен как классификация препятствий}.
\end{enumerate}
\end{protocol}

Следующей проверки внутри Тома V нет. Новая версия программы может быть
открыта только отдельным решением с явной заменой аксиоматики и новым
критерием завершённости.

Машинный реестр сохранён в
\path{s2t_v5_parent_architecture_status_freeze_gate_results.json}.